\chapter{Discussion and System Synthesis}
\label{ch:synthesis}

This chapter steps back from the two systems to their joint meaning. \Cref{sec:discussion} synthesizes
what intent-conditioned single-path selection and joint multipath co-optimization say together about
intent-driven network control. \Cref{sec:analysis} then analyzes both systems along the axes that decide
whether learned path control is practical -- overhead, scalability, and generalization -- and
\cref{sec:analysis:deployability} states the deployment gaps and trust assumptions under which the
controllers operate. Where a claim depends on measurements still in progress, it is marked.

\section{Synthesis of Findings}
\label{sec:discussion}

Read together, the two systems tell a consistent story about \emph{when} additional control machinery
earns its keep.

\subsection{Interpretation}
\label{sec:disc:interpretation}

\paragraph{Intent conditioning turns a family of policies into one.} The central lesson of
\cref{ch:single_path} is that the right way to serve many objectives is not many models but one model that
takes the objective as input -- provided the objective reaches the part of the network that actually
decides. \Cref{prop:p1cancel} is the crisp version of this point: a conditioning signal that enters
identically across the alternatives cancels under the comparison that picks the action, so the policy can
look conditioned while behaving identically. Supplying the objective to the per-path terms avoids this,
because the objective can then interact with each candidate's own feature values and change which one
wins. This is a design principle, not a SCION-specific trick: any comparison-based selector conditioned
on an objective must route that objective into its per-alternative terms. It holds in both settings this
thesis studies, and \cref{prop:p2cancel} makes it sharper in the second --- under a softmax allocation
the cancellation is exact on the entire action, not merely on which alternative wins.

\paragraph{But the guarantee is narrower than the behavior.} The multipath system supplies the
qualification, and it is the more useful half of the result. \Cref{prop:p2cancel}'s premise is that the
conditioning be additive \emph{on the terms the comparison consumes}, and three independent mechanisms
in a perfectly ordinary architecture break it: a saturating nonlinearity between the conditioning and
the comparison, a shared encoder that mixes the conditioning into each alternative's hidden
representation, and an observation normalizer that is itself a function of the objective. Measured on
trained policies, a globally conditioned agent that theory predicts should be inert is instead a working
one, indistinguishable in behavioral divergence from the per-alternative variant
(\cref{sec:p2eval:conditioning}). The reconciliation is that the theory says where conditioning must
enter to be \emph{guaranteed} to act, while a deep encoder learns to get there without a guarantee. Both
statements are needed. A designer who has only the first over-engineers; one who has only the second is
relying on an accident of the function class, and \cref{ch:single_path}'s value-stream agent --- which
did re-rank three times less --- is what that accident looks like when it does not happen.

The third of those three mechanisms turned out to carry a measurable cost, which is the one place in
this thesis where a premise violation is not merely a caveat on a proof. An observation normalizer that
is itself a function of the objective sends the objective into each alternative's features by a second,
unintended route, and under per-path conditioning the policy then has to reconcile two correlated
channels where the design specifies one. Closing that route recovers $+0.036$ and $+0.033$ QoE at the
two intents that had a deficit, changes nothing at the one that did not, and tightens the training-seed
spread (\cref{sec:p2eval:conditioning}). The design lesson is narrow and portable: when conditioning is
routed into per-alternative terms, audit the \emph{preprocessing} for the same signal, because a
normalizer derived from the objective is a conditioning path that nobody declared.

\paragraph{The two decisions are not separable.} The first lesson of \cref{ch:multipath} is the one the
design was built to test: removing either agent forfeits most of the achievable QoE. An application
agent over a naive split recovers as little as 35\% of the joint system, and a learned split under a
conventional bitrate controller as little as 84\% (\cref{sec:p2eval:ablation}). Those recovery ratios
hold only inside the regime where the system works: on a harsher held-out network both terms go negative
and the ratio stops meaning anything, which is why \cref{sec:p2eval:regimes} reports the absolute
difference instead. Rate and split are
coupled through the same bottleneck, and a controller that fixes either one is optimizing against a
constraint it could have moved. The decomposition is a claim about the analytical backend: on the
packet-level one the application-only ablation recovers 80\% rather than 35\%, and the network-only
ablation costs nothing measurable at all. What the packet-level evidence supports is that the learned
\emph{scheduler} is worth having; that the learned bitrate controller on top of it is worth having
remains a mock-backend result.

\paragraph{What governs the payoff is headroom, not motion.} The second lesson is a correction to the
intuition we began with. It is tempting to say that learned multipath scheduling matters when the
network is volatile, and the correction matters because that framing would misdirect anyone deciding
whether to deploy it. The governing quantity is the relationship between the load an even split places
on each path and that path's residual capacity (\cref{sec:p2eval:regimes}). Churn, regime shifts, and
bursts matter because they destroy that evenness transiently --- but they are not the variable. A
stationary but \emph{asymmetric} topology defeats application-only adaptation as thoroughly as a
churning one, while a stationary topology of near-equal paths leaves the scheduler almost nothing to
earn. A third way of destroying the same evenness makes the point without touching the dynamics at all:
dividing a \emph{fixed} aggregate capacity among more paths raises the scheduler's contribution
monotonically, from nothing at two paths, where an even split is already near-optimal, to the point at
twelve where disabling it drives QoE negative (\cref{sec:p2eval:pathcount}).

Three practical implications follow. One should not benchmark a multipath scheduler on a
well-balanced topology, stationary or not, because that is the regime in which any allocation works ---
and, as \cref{sec:p2eval:regimes} shows, the regime in which the measured margin can fall below the
spread between training runs of the same method, so the benchmark reports its own seed draw. The
payoff scales with the application's intent as well as with the network: the stricter the delay budget,
the less slack there is to absorb a poor allocation, so the more the allocation matters. And the two
axes are \emph{not} separable, which the two-dimensional map of \cref{sec:p2eval:regimes} is what
establishes: the three intents differ by 0.039 QoE on a benign network and by 0.566 on a harsh one, a
factor of 14.5, so reading the two marginals additively understates the interaction badly. A
practitioner who measures the intent axis on a well-behaved network will conclude that intent hardly
matters, and will be measuring the network rather than the intent.

\paragraph{A surrogate can be good enough to train in and useless to choose in.} The last lesson is
methodological and was not one we set out to learn. Six policies trained identically on the analytical
backend are indistinguishable there --- a spread of 0.008, smaller than a single arm's evaluation-seed
interval --- and separate into two clean clusters on the packet-level backend, three retaining their
first-place ranking and three not. The rank correlation between the two backends is $-0.086$: absent, not
weak, with the best surrogate checkpoint the worst on the realistic body
(\cref{sec:p2eval:ns3}). The mechanism is identifiable and mundane. The variable that decides the
packet-level outcome is delay-weighted allocation --- where the bytes go, priced by what each path costs
at that moment --- and the surrogate's analytical queue compresses it into a \SI{1.9}{\milli\second} band
where the realistic body spans \SI{9.7}{\milli\second}. A model can therefore be faithful enough to learn
a good policy in and blind to the dimension that separates good policies from bad ones. The practical
rule is narrow and firm: select checkpoints on the backend the claim is about, or report the
distribution over training seeds and select nothing.

\paragraph{Conditioning pays operationally only when the intent moves.} A result we did not anticipate
bounds what the conditioning argument is worth in deployment. With the objective fixed for the duration
of a flow -- the regime every headline number in this thesis is measured in -- one conditioned policy
and a family of specialists are operationally interchangeable, because a deployment can simply load the
right specialist when the flow starts. The advantage is real only when the objective changes
\emph{while the connection is up}, and there it is one-sided: relaxing a deadline mid-call costs a
policy that cannot see the change nothing measurable, while tightening one drives a specialist six
times past its new budget for ten seconds where the conditioned policy never crosses it
(\cref{sec:p2eval:conditioning}). The practical statement is therefore narrower and more useful than
``one model instead of many'': the reason to condition is that objectives move, and the cost of not
conditioning is paid entirely in the direction that tightens the constraint.

\paragraph{One principle, two layers.} Both systems instantiate the Mortise principle~\cite{Shen2026} --
intent as a first-class input that reshapes a network mechanism -- at successively deeper points of
control. The single-path selector reshapes \emph{which path} is chosen; the multipath system reshapes both
\emph{how the paths are used} and \emph{what the application produces}. In both, one conditioned policy
replaces a family of per-objective ones. The progression suggests a general recipe: expose the
application's intent, and let it steer every controllable decision between the application and the wire.

\subsection{Implications}
\label{sec:disc:implications}

The results matter for how we build controllers on path-aware architectures. Prior learned path
selectors~\cite{Bourak2025} and learned congestion or scheduling controllers~\cite{Jay2019, Zhang2019,
Wu2020} each optimize a fixed objective; our work argues that on a path-aware substrate the objective should
be a knob, and that a single conditioned policy can replace a zoo of per-objective models. For real-time
media specifically, the joint bitrate-and-split result reinforces a theme from learned
ABR~\cite{Mao2017} -- that application and network decisions are better made together -- and extends it to
the multipath setting that SCION makes routine. More broadly, the two systems are evidence that SCION's
endpoint path choice is not merely a routing convenience but an enabler of application-aware optimization
that is awkward or impossible in a destination-based Internet.

\subsection{Limitations}
\label{sec:disc:limitations}

The results are bounded in ways that the analysis (\cref{sec:analysis}) details and that fairness requires
restating here. Both systems are evaluated in simulation. The multipath system's headline numbers are
measured on the analytical mock, and on the packet-level backend they hold for three of six independently
trained policies and not for the other three, with absolute QoE not portable in any case
(\cref{sec:p2eval:ns3}). Every reported policy is trained on the mock; two packet-level-native training
seeds do learn, but two is below the three-seed floor the chapter applies elsewhere, so whether the
system \emph{learns} on the realistic body is supported rather than established
(\cref{sec:p2impl:gaps}). Since mock score does not predict packet-level rank, the mock's standing as a
training surrogate is an open question rather than a technicality. The headline multipath result is also
bounded to the distribution it was measured on: it survives on a milder held-out network and fails on a
harsher one, where a hand-written GCC scheduler beats it in every column
(\cref{sec:p2eval:regimes}). The parity of the intent-conditioned policy with the per-intent specialists is a failure to
detect a difference, not a demonstration of equivalence, and we can now say how weak a failure: at six
training seeds per arm the design resolves nothing below \SIrange{0.063}{0.087}{} QoE, while every
margin at issue is 0.048 or less, so the experiment could not have separated the two outcomes either
way (\cref{sec:p2eval:conditioning}). The specialists it matches were themselves not measurably
specialized, for the same reason (\cref{sec:p2eval:transfer}).
The single-path study rests on one topology and one training seed, with a purely temporal train/test
split, so its quantitative margins should be read as indicative rather than as established effect sizes
(\cref{sec:p1eval:setup,sec:p1eval:ablation}). Intent in this thesis is \emph{declared}, not
\emph{inferred}: we assume the application can state its objective, and we do not study how to recover a
latent intent from traffic -- a natural but separate problem. The transport model is TCP-like rather than true Multipath QUIC, and the VMAF surrogate is
bitrate-dominated, so absolute QoE figures should be read as relative comparisons rather than calibrated
user scores. Finally, we have not addressed online, safe learning against live traffic, which any
deployment would require. None of these undercut the two qualitative findings -- conditioning that reaches
the per-alternative comparison is the only kind guaranteed to act, and joint rate-and-split control pays
off in proportion to how unevenly path capacity is distributed and how strictly the application prices
delay -- but they do delimit how far the quantitative results should be pushed.

\gap{two limitations are missing from a list that is otherwise exemplary, and both are ones a reader
will find on their own if this section does not state them. (i) \emph{\Cref{ch:multipath} does not run
on SCION.} There is no control plane, no beaconing, no path store; the six paths are a hand-specified
configuration (\cref{tab:app:p2topo}) and the study assumes only path multiplicity. Given the thesis
title, this belongs in the limitations list explicitly rather than being inferable from
\cref{sec:mp:impl}. (ii) \emph{\Cref{ch:single_path}'s seed coverage is partial.} The paragraph above
says ``one topology and one training seed'', which understates the situation in one direction and
overstates it in another: the ablation ladder \emph{is} at five seeds with disjoint intervals
(\cref{tab:p1eval:seeds}), while the intent matrix, zero-shot sweep, path-count study, probing
comparison and ceiling figure are all single-run. Say which is which. The asymmetry matters because
\cref{sec:p2eval:seeds} makes single-run reporting the thesis's own cautionary tale, and a limitations
section that blurs the distinction invites the reader to discount the results that \emph{are} properly
seeded along with those that are not.}

\section{System Analysis}
\label{sec:analysis}

This section analyzes the two systems along the dimensions that determine whether learned path control is
practical: overhead (\cref{sec:analysis:overhead}), scalability (\cref{sec:analysis:scalability}), and
generalization (\cref{sec:analysis:generalization}).

\subsection{Overhead}
\label{sec:analysis:overhead}

Two kinds of overhead matter: measurement and computation.

\paragraph{Measurement (probing) overhead.} In the single-path selector (\cref{ch:single_path}), probing
is a first-class cost charged in the reward (\cref{eq:p1reward}) and accounted per decision
(\cref{sec:p1impl:baselines}). The learned selector probes only the paths it inspects, in contrast to
heuristics that require a metric on every path, so its measurement footprint is a fraction of an exhaustive
scan at comparable selection quality -- and, in fact, smaller than that of every heuristic baseline, since
each of those must measure all candidates before applying its rule (\cref{sec:p1eval:overhead}). In the
multipath system (\cref{ch:multipath}), per-path state is derived from ordinary transport feedback (RTT
samples, acknowledgements, loss signals) rather than dedicated probes, so the marginal measurement cost of
the split decision is low.

\paragraph{Computational overhead.} Both policies are small neural networks evaluated once per decision.
The single-path selector runs at path-selection granularity (seconds to minutes) and is comfortably cheap.
The multipath transport agent is the tighter constraint: it runs \emph{per frame}
(${\approx}\SI{33}{\milli\second}$ at 30\,fps), so its inference must fit well inside a frame interval. The
scoring architecture keeps this affordable -- a shared per-path encoder evaluated $N$ times and a softmax --
but the per-frame budget is the binding real-time constraint. For the single-path selector the margin is
ample: timing only the decision itself, over 1280~held-out contexts of 30~candidate paths each on a
22-core CPU host with no GPU, the conditioned selector decides in under \SI{0.3}{\milli\second} on
average and under \SI{2}{\milli\second} at the 99th percentile across repeated runs. That is three orders
of magnitude inside a path-selection budget of seconds to minutes, and comfortably inside a
\SI{33}{\milli\second} video frame. The multipath transport agent, whose per-frame budget is the
tighter of the two, has the same shape of answer: it decides in \SI{0.507}{\milli\second} on average
on a 16-core CPU host with no GPU, about 1.5\% of the frame interval it must fit inside, and
\SI{0.71}{\milli\second} at the 99th percentile over the 40{,}500 decisions of a full nine-cell grading
(\cref{fig:p2eval:decision,sec:p2eval:pathcount}). Both controllers are therefore compute-bound nowhere
near their deadlines. The cost grows linearly in the candidate count in each case, since the shared
per-path encoder is evaluated once per path, but for the multipath agent the slope is what makes the
margin large rather than the intercept: measured over $N \in [2, 96]$ the fit is
$t(N) = 0.483 + 0.0050\,N$\,ms at $R^2 = 0.990$, so 94\% of the cost at six paths is a constant and each
additional path costs \SI{5}{\micro\second}. The frame budget binds near 6500 paths.

One qualification belongs with any real-time claim and is easy to lose behind a mean. The 99.9th
percentile is still comfortable at under \SI{3.6}{\milli\second}, but roughly 0.01\% of decisions
overrun \SI{33.3}{\milli\second} outright, with a worst case of \SI{170}{\milli\second} to
\SI{219}{\milli\second}. That rate is independent of the path count and appears at two paths as readily
as at twelve, which identifies it as a Python-runtime artifact --- allocation and garbage collection ---
rather than model cost. It would not be fixed by a smaller network; it would be fixed by a pre-allocated
inference path, which any deployment would want regardless.

\subsection{Scalability}
\label{sec:analysis:scalability}

The scoring architecture is what makes both systems scale in the number of paths. Because a shared network
scores each path independently and decisions are made by comparison (\cref{sec:p1design:scoring}) or softmax
(\cref{sec:p2design:scoring}), the parameter count is independent of $N$, and inference is linear in $N$.
Both halves of that are now measured rather than asserted, in both systems. The multipath transport
policy is 69,122 parameters at every path count from 2 to 96, and one six-path checkpoint runs
unmodified at 2, 3, 4, 6, 8, and 12 live paths, ranking first in all nine evaluation cells at every count
from four upward, while the flat agent it replaces cannot even be loaded at a path count other than its
own (\cref{sec:p2eval:pathcount}). The single-path selector holds regret two orders of magnitude below
the fixed-action agent across a 15-fold range in candidate-set size (\cref{sec:p1eval:dynamic}). One
trained model therefore serves source--destination pairs with any number of paths, without retraining or
padding to a fixed maximum, which is the limitation of the flat agents.

Two limits qualify this. Path-count invariance is not capacity invariance: the same multipath policy
collapses when the aggregate bandwidth behind those paths falls to a third of what it trained on, and
ranks first in zero of nine cells there, though it is untroubled above the training point
(\cref{sec:p2eval:pathcount}). And the remaining scalability question is training coverage --- the
environment must expose enough variety of path counts, topologies, and conditions for the shared scorer
to generalize --- which the path-count sweep does not answer, because it measures \emph{transfer} from a
six-path checkpoint rather than what the head would achieve if trained at each count.
\gap{report how selection quality holds as topology size grows, ideally on topologies larger than
those seen in training. The path-count axis is now measured in both systems; the topology-size axis is
not, in either.}

\subsection{Generalization}
\label{sec:analysis:generalization}

Generalization is required along three axes. \textbf{Across intents:} the conditioning of
\cref{ch:single_path} is meant to interpolate to intents between (and, ideally, beyond) the trained reward
profiles, so that a novel weight vector yields sensible behavior without retraining.
\Cref{sec:p1eval:zeroshot} confirms it: sweeping the intent along the segment between two trained
profiles moves the chosen path's latency down smoothly and, in four of five training seeds, strictly
monotonically (Spearman $-1.000$; the fifth seed has one ascending step of \SI{0.036}{\milli\second}, for
$-0.991$), with no seed carrying more than 31\% of the span in a single step, and modest extrapolation
beyond both endpoints continues the trend. The multipath conditioning generalizes along the same axis but with a weaker
warrant: sweeping the intent between two profiles moves bitrate and quality smoothly and monotonically,
by 11--24 times the span a non-learning control shows over the same rollouts, but training there sampled
the intent space continuously, so the interior of the segment is in-distribution rather than unseen
(\cref{sec:p2eval:conditioning}). At six training seeds that claim also has to be read per seed rather
than pooled: the per-path arm's trend is positive in all six runs, the global-only arm's in four of six,
with two trending the wrong way behind a pooled statistic that looks clean. Extrapolation past the endpoints fails asymmetrically in both systems,
and in the multipath case informatively so: loosening the intent beyond anything trained holds QoE,
tightening it beyond the trained range collapses it -- which is the expected direction, since the latter
asks for a deadline below what the network can deliver rather than merely an unfamiliar one.
\textbf{Across topologies and path counts:} the permutation-equivariant design gives
structural generalization, and \cref{sec:p1eval:dynamic} confirms both halves of it empirically -- exact
order invariance over all 483840~permutation trials, across three agents and five training seeds with
zero mismatches, and regret held two orders of magnitude below the fixed-action agent across a 15-fold
range in candidate count. \Cref{sec:p1eval:ablation} adds a
structurally disjoint set of 32~pairs spanning 32~source ASes never visited in training, on which the
conditioned selector's optimality gap grows only from 0.0018 to 0.0033 while the flat agent's Throughput
reward goes negative. What remains untested is scale: all of this is one topology of 90~ASes.
\textbf{Across conditions:} the dynamic scenario of
\cref{ch:multipath}
(\cref{sec:p2eval:results}) is itself a generalization stress test, since the policy must cope with churn
and bursts rather than a stationary environment; the collapse of the fixed-schedule baselines there shows
that a fixed allocation rule does \emph{not} survive a network whose headroom keeps shifting underneath
it, which is the gap learned control fills. How far that extends beyond the training distribution is now
mapped, and the answer is asymmetric: graded over five network regimes spanning a milder held-out
configuration to a harsher one, the learned pair keeps first place in every cell at every intent on the
mild and mild--nominal rows, and loses it entirely on the harsh row, where it ranks first in zero of nine
cells at two of three intents and scores below plain GCC (\cref{sec:p2eval:regimes}). Its training-seed
spread rises with harshness from 0.018--0.024 to 0.051--0.339, so out of distribution the learner is not
merely worse but inconsistent. The headline multipath claim should therefore be quoted with the
distribution it was measured on --- it generalizes downward in difficulty and not upward.
For the single-path system, off-distribution degradation is bounded on two of these three axes:
unseen intents cost nothing measurable (\cref{sec:p1eval:zeroshot}), and unseen pairs and candidate-set
sizes cost at most about 0.01~reward (\cref{sec:p1eval:dynamic,sec:p1eval:ablation}).
The held-out-dynamics axis is now measured for the multipath system and the answer is a clean negative:
graded on two configurations whose hazard rates lie outside anything either arm trained over, domain
randomization costs a further 0.033 QoE in nine of nine cells, on top of the 0.027 it costs on the
nominal environment (\cref{sec:p2eval:ablation}). The more permissive reading that was available before
this run -- that randomization trades nominal performance for robustness -- is not supported: it loses on
both axes here. One of the two held-out environments is uninformative rather than contrary, both arms
scoring negative under rankings that do not hold.

\gap{what remains unquantified on this axis is scale, for both systems: the single-path results are one
topology of 90~ASes and the multipath results one six-path configuration, so nothing here bounds decay
as topology size grows. This is the more expensive of the two remaining generalization questions and the
only one still open.}

\section{Deployment Considerations and Trust Assumptions}
\label{sec:analysis:deployability}

\subsection{Deployment Gaps}
Several gaps still separate these prototypes from a deployment.

\begin{itemize}
  \item \textbf{From simulation to reality.} Both systems are trained in simulation. The single-path study
  relies on a synthetic traffic and beaconing environment; the multipath system on NS-3 and a mock
  transport. Bridging the gap to deployment would require sim-to-real transfer -- for example through domain
  randomization or online fine-tuning against live measurements -- and the packet-level NS-3 backend is the
  intended stepping stone (\cref{sec:p2impl:gaps}).
  \item \textbf{Real transport.} The multipath system models subflows with TCP-like congestion control
  rather than true Multipath QUIC~\cite{DeConinck2017}. A real QUIC integration is therefore needed to
  validate the per-frame scheduling under an actual multipath transport.
  \item \textbf{Online and safe learning.} A deployed agent would learn or adapt against live traffic, where
  exploration has real cost. Bounding worst-case behavior -- for example by falling back to a safe heuristic
  when the policy is uncertain, or by constraining exploration -- is a prerequisite that this work does not
  address.
\end{itemize}

Two properties help on the positive side. First, the endpoint-driven nature of SCION path selection means
these controllers require no network-internal changes. Second, the horizon-agnostic application agent
(\cref{sec:p2design:hierarchy}) is designed for unbounded real sessions rather than fixed-length episodes.

\subsection{Trust Assumptions}

This thesis is a performance and control study, not a security study, so in place of a classical
attacker model we state the trust assumptions under which the learned controllers operate and note
where security considerations would enter a real deployment.

\begin{itemize}
  \item \textbf{Honest measurements.} We assume that the per-path measurements an agent acts on --
  whether obtained by active probing (\cref{ch:single_path}) or transport feedback (\cref{ch:multipath}) --
  are honest. A malicious on-path element that fabricates favorable metrics could steer the agent toward a
  bad path; defending against that is out of scope, although the trust term below offers partial mitigation.
  \item \textbf{SCION path authenticity and trust.} SCION's cryptographically authenticated paths mean
  an agent knows \emph{which} ASes a path traverses and cannot be silently redirected. The single-path
  reward includes a \emph{trust} component derived from path quality, allowing the agent to prefer
  higher-confidence paths and discount those with anomalous loss or latency -- a lightweight hedge against
  unreliable or adversarial path advertisement, not a full defense.
  \item \textbf{No adversarial control of the learner.} We assume the agent's own training and
  inference pipeline is trusted; adversarial manipulation of the RL policy (e.g.\ reward poisoning or
  observation attacks) is beyond this work.
\end{itemize}

A production deployment would need to revisit these assumptions -- in particular by verifying measurement
integrity and by bounding the damage a single misreported path can do.
